% !TEX root = lectures.tex
\newcommand{\Exp}{\mathbb{E}}
\section{Lecture 2}
I was out.

\section{Lecture 3}
\subsection{Concentration of Variance}
We will show a subgaussian decay tail bound,
by showing we can convert from the expectation versions of the bounds to
high probability ones.

Note the standard fact $\Var(\sum_i X_i) = \sum_i \Var{X_i}$
for independent $X_i$'s. Also $\Var(x) = \frac{1}{2} \Exp_{x, y} (x - y)^2$
for $(x, y)$ as two independent copies of $y$.

\begin{theorem}[Efron-Stein]
    Suppose $X_1, \dots, X_n$ are independent, real-valued random variables.
    Then for all $f: \R^n \to \R$,
    \[ \Var(f(X_1, \dots, X_n)) \le \sum_{i = 1}^{n} \mathbb{E}_{x \sim i} \Var_{x_i} f(x) \]
    Where we mean
    \begin{align*}
        \mathbb{E}_{x \sim i} \Var_{x_i} f(x) &:= \mathbb{E}_{x_1, \dots, x_{i - 1}, x_{i + 1}, \dots, x_n} \Exp_{x_i} (f(x_1, \dots, x_n) - \Exp_{x_i} f(x_1, \dots, x_n))^2 \\
        &= \frac{1}{2} \Exp_{x, y} (f(x_1, \dots, x_{i - 1}, x_i, x_{i + 1}, \dots, x_n) - f(x_1, \dots, x_{i - 1}, y_i, x_{i + 1}, \dots, x_n))^2
    \end{align*}

    \begin{proof}
        Define $(\Exp_i f)(x_1, \dots, x_{i - 1}, \cdot, x_{i + 1}, \dots, x_n) := \Exp_{x_i} f(x_1, \dots, x_{i - 1}, x_i, x_{i + 1}, \dots, x_n)$.
        Define $\Exp_{S} f = \Exp_{i_1} \Exp_{i_2} \dots \Exp_{i_{|S|}} f$.
        Notice that $f = (f - \Exp_{[1]} f) + (\Exp_{[1]} f - \Exp_{[2]} f) + \dots + (\Exp_{[n - 1]} f - \Exp_{[n]} f) + \Exp_{[n]} f$.
        Notice that $\Exp_{[n]} f = \Exp f$. Then, we claim that in the quantity $\Exp(f - \Exp f)^2$, the cross terms cancel.
        For example, 
        \[ \Exp [(f - \Exp_{[1]} f)(\Exp_{[1]} f - \Exp_{[2]} f)] = \Exp f \Exp_{[1]} f - \Exp f \Exp_{[2]} f - \Exp (\Exp_{[1]} f)^2 + \Exp(\Exp_{[1]} f)(\Exp_{[2]} f) = 0\]
        where the first two terms cancel by the tower rule, and same with the second two. Now,
        $\Exp(f - \Exp f)^2 = \sum_{i = 0}^{n - 1} \Exp(\Exp_{[i]} f - \Exp_{[i + 1]} f)^2$.
        Now, we can apply Jensen's inequality term by term, to get an expecation over the correct variables.
    \end{proof}
\end{theorem}


\begin{definition}
    Consider $\Sigma \subseteq \Sigma$ and $f: \Sigma^n \to \R$.
    \begin{align*}
        D_i f(x) &= \sup_{y \in \Sigma} f(x_1, \dots, x_{i - 1}, y, x_{i + 1}, \dots, x_n) - \inf_{y \in \Sigma} f(x_1, \dots, x_{i - 1}, y, x_{i + 1}, \dots, x_n)\\
        \Delta_i f &:= \sup_{x \in \Sigma^n} (D_i f(x))
    \end{align*}
\end{definition}
Note that if $x \in [a, b]$ almost surely, then, $\Var(x) \le \frac{1}{4} (b - a)^2$.
\begin{corollary}
    \[ \Var(f(x_1, \dots, x_n)) \le \frac{1}{4} \sum_{i = 1}^n \Exp (D_i f(x))^2 \le \frac{1}{4} \sum_i (\Delta_i f)^2 \]
\end{corollary}

\begin{theorem}
    Let $W \in \R^{n \times n}$ be symmetric, independent entries such that $|W_{ij}| \le K$.
    Then, $\Var(\lambda_1(W)) \le 16 K^2$.
\end{theorem}


\begin{corollary}
    Any Erdos-Renyi random graph's largest eigenvalue has variance at most $16$.
\end{corollary}

\begin{definition}
    $D_i^+ f(x) = \sup_{y \in \Sigma} f(x_1, \dots, x_{i - 1}, y, x_{i + 1}, \dots, x_n) - f(x)$.
\end{definition}

Then the corollary before turns into.
\[ \Var(f(x_1, \dots, x_n)) \le \frac{1}{4} \sum_{i = 1}^n \Exp (D_i^+ f(x))^2 \le \frac{1}{4} \sum_i (\Delta_i^+ f)^2 \]
Now we can prove the variance theorem.
\begin{proof}
    We now will control $D_{ij}^+ \lambda_1(W)$.
    Suppose that $W \neq W'$ differ in only the $(i, j)$th entry. Let $v$
    be the top eigenvector of $W$. 
    \[ \lambda_1(W) - \lambda_1(W') = v^T W v - \max_{\norm{u} = 1} u^T W' u = v^T W v - v^T W' v \le 2 |(W_{ij} - W_{ij}')| |v_i v_j| \le 4K |v_i v_j|\]
    Then by the new corollary we have $\Var(f) \le 16K^2 \sum_{i, j} v_i^2 v_j^2 = 16K^2$.
\end{proof}

\begin{theorem}[Gaussian Poincare Inequality]
    Let $f: \R^n \to \R$ be $\mathcal{C}^1$ and
    \[ \Var_{x \sim \N(0, I)} f(x) \le \mathbb{E} \norm{(\nabla f)(x)}_2^2 \]
    
    \textbf{Proof.}
        By Efron-Stein we have that
        \begin{align*}
            \Var_{x \sim \N(0, I)} (f(x)) &\le \sum_{i = 1}^n \Exp_{x \sim i} \Var_{x_i} f(x) \\
            &\le \sum_{i = 1}^{n} \Exp_{x \sim i} \Exp_{x_i} \qty(\partial_i f(x))^2 \\
            &= \Exp_{ x \sim \N(0, I)} \norm{\nabla f(x)}_2^2
        \end{align*}
        if we assume the thorem for $n = 1$.

        Now, notice 
        \begin{align*}
            \Var_x f(x) &\approx_{CLT} \Var_{z} f\qty( \frac{1}{\sqrt{k}} \sum_{i = 1}^k z_i) \\
            &\le \frac{1}{2} \sum_{i = 1}^{k} \Exp_{z, z'} f\qty(\frac{1}{\sqrt{k}} \sum_{j \neq i} z_j + \frac{z_i}{\sqrt{k} }) \\
            &= \frac{1}{4} \sum_{i = 1}^{k} \Exp_{z} \qty(f\qty(\dots + \frac{1}{\sqrt{k}}) - f\qty(\dots - \frac{1}{\sqrt{k}}))^2 \\
            &\approx_{taylor} \frac{1}{4} \sum_{i = 1}^k \Exp_z \qty(\frac{2}{\sqrt{k}} f'\qty(\frac{1}{\sqrt{k}} \sum_{j \neq i})) \\
            &= \frac{1}{k} \sum_{i = 1}^{k} \frac{1}{4} \Exp_z \qty(\frac{2}{\sqrt{k}} f'\qty(\frac{1}{\sqrt{k}} \sum_{j \neq i} z_j)) \\
            &\approx_{CLT} \Exp f'(x)^2
        \end{align*}
\end{theorem}

\begin{corollary}
    If $f: \R^n \to \R$ is $L$-Lipschitz,
    $|f(x) -f(y)| \le L \norm{x - y}_2$
    \[ \Var_x f(x) \le L^2 \]
\end{corollary}

\begin{theorem}
    Let $W \in \R^{n \times n}$ be a symmetric random matrix with independent entries and assume $W_{ij} = W_{ji} \sim \N(\mu_{ij} ,\sigma_{ij}^2)$.
    Then,
    \[ \Var(\lambda_1(W)) \le 2 \max_{i,j} \sigma_{ij}^2 \]
\end{theorem}

\subsection{Subgaussian Tail Bounds}
\begin{theorem}[Gaussian Concentration]
    Let $x \sim \N(0, I)$, $F: \R^n \to \R$ be $\sigma$-Lipschitz (in 2-norm),
    then for all $t \ge 0$, $\mathbb{P}_x[|F(x) - \E{f}| \ge t] \le 2e^{-\frac{t^2}{2\sigma^2}}$.

    \begin{proof}
        We will prove a slightly weaker bound of $2e^{-\frac{t^2}{\pi^2 \sigma^2}}$.
        Assume $F$ is $\mathcal{C}^1$ (otherwise we can "smoothen" it out). Then, for all $x$,
        $\norm{\nabla F (x)}_2 \le \sigma$. Without loss of generality, $\E{F} = 0$,
        otherwise, we will use $F - \E{F}$.
        We will prove the one-sided bound; the other side follows by using the function $-F$.

        Now, we can use the exponential moment method to write
        \[ \Pr [F(x) \ge t] = \Pr[\exp(\lambda F(x)) \ge \exp(\lambda t)] \le \frac{\Exp[\exp(\lambda F)]}{\exp(\lambda t)} \]
        Let $y$ be an independent copy of $x$.
        \[ \Exp[\exp(-\lambda F(y))] \ge \exp(-\Exp \lambda F(y)) = \exp(0) = 1 \]
        Then,
        \[ \Exp[\exp(\lambda(F(x) - F(y)))] = \Exp[\exp(\lambda F(x))] \Exp[\exp(-\lambda F(y))] \ge \Exp[\exp(\lambda F)] \]
        The key trick is that
        \[ F(x) - F(y) = \int_{0}^{\pi/2} \pdv{}{\theta} F(y \cos \theta + x \sin \theta) \dd{\theta} \]
        Calling $x_{\theta} = y\cos \theta + x \sin \theta$, note that $x_{\theta} \sim \N(0, I)$.
        Then, $x_{\theta}' = \pdv{x_{\theta}}{\theta} = -y \sin \theta + x \cos \theta \sim \N(0, I)$.
        The two combinations are orthogonal hence $x_{\theta}$ and $x_{\theta}'$ are independent.
        Thus,
        \begin{align*}
            \exp(\lambda(F(x) - F(y))) &= \exp(\lambda \frac{\pi}{2} \frac{1}{\pi/2} \int_0^{\pi/2} \pdv{}{\theta} F(x_{\theta}) \dd{\theta} ) \\
            &= \exp(\lambda \frac{\pi}{2} \Exp_{\theta \in [0, \pi/2 ]}[\pdv{}{\theta} F(x_{\theta})] ) \\
            &\le \Exp_{\theta \in [0, \pi/2]} \exp(\lambda \frac{\pi}{2} \pdv{}{\theta} F(x_{\theta})) \\
            &= \Exp_{x,y} \frac{2}{\pi} \int_{0}^{\pi/2} \exp\qty(\lambda \frac{\pi}{2} \langle \nabla F(x_{\theta}), x_{\theta}' \rangle) \dd{\theta} \\
        \end{align*}
        Notice that $\langle \nabla F(x_0), x_{\theta}' \rangle$ is distributed as a Gaussian
        with variance at most $\sigma^2$, since the gradient is at most $\sigma$ in magnitude and the two vectors are independent.
        \begin{align*}
            \exp(\lambda(F(x) - F(y))) &\le \frac{2}{\pi} \int_0^{\pi/2} \Exp_{x, y} \exp\qty(\frac{1}{2} \qty(\lambda \frac{\pi}{2} \sigma)^2) \dd{\theta} \\
            &\le \exp(1/2(\lambda \pi/2 \sigma)^2)
        \end{align*}
        which implies our bound by choosing $\lambda = \qty(\frac{2}{\pi}) \frac{t}{\sigma^2}$.
    \end{proof}
\end{theorem}

\begin{corollary}
    Let $X = \sum_i g_i A_i$, $g_i \sim \N(0, 1)$ independently, $A_i \in \R^{n \times n}$, symmetric,
    $\lambda_1(X)$ and $\norm{X}_2 = \sigma_1(X)$ are both $\sigma_*(x)^2$-subgaussian,
    where
    \[ \sigma_*(x)^2 = \max_{\norm{v}_2 = 1} \langle \sum_{i = 1}^m A_i^{\otimes 2}, v^{\otimes 4} \rangle = \max_{\norm{v}_2} \sum_{i = 1}^{m} (v^T A_i v)^2 \]
    \begin{proof}
        Let's say we want to prove this for $\lambda_1$. Then 
        \[\lambda_1(X(g)) - \lambda_1(X(h)) \le \norm{X(g) - X(h)}_2 =  \norm{\sum_{i = 1}^m(g_i - h_i) A_i}_2.\]
        Then, 
        \[\sum_{i = 1}^{m} (g_i - h_i) v^T A_i v \le \qty(\max_{\norm{v}_2 = 1} \sum_{i = 1}^{m} (v^T A_i v)^{2})^{1/2} \norm{g - h}_2\]
        and we can apply the previous theorem with this Lipschitz constant.
        (The spectral norm version follows from triangle inequality.)
    \end{proof}
\end{corollary} 

Recall that previously, the error term was
\[ \Exp_g \norm{\sum_i g_i A_i}_2 \le C \sqrt{\log n} \norm{\sum_i A_i^2}_2^{1/2} \]
Then, $\sigma(x)^2 = \max_{\norm{v}_2 = 1} \sum_i v^T (A_i^2) v$.
Now, our 
\[ \sigma_*(x)^2 = \sum_{i = 1}^{m} \langle A_i v, v \rangle^2 \le \sum_{i = 1}^{m} \norm{A_1 v}_2^2 \norm{v}_2^2 = \sum_{i = 1}^{m} v^T A_i^2 v = \sigma^2 (x) \]

\begin{corollary}
    Let $W \in \R^{n \times n}$, $W_{ij} = W_{ji} \sim \N(0, 1)$, $W_{ii} \sim \N(0, 2)$ (Wigner distribution).
    Then,
    \[ \Pr{\norm{W}_2 \ge \Exp \norm{W} + t} \le 2e^{-t^2/4} \]
\end{corollary}

\begin{theorem}[Talagrand Concentration]
    Let $|X_i| \le k$ be arbitrary bounded independent random variables, let $F: \R \to \R$, $\sigma$-Lipschitz, convex.
    then for all $t \ge 0$
    \[ \Pr_x[|F(x) - \E{f}| \ge tK] \le 2e^{-t^2/2\sigma^2} \]
\end{theorem}

This will give tail bounds for $\norm{X}_2$ with Rademachers, but not eigenvalues. Convexity is required,
even over the hypercube.

\begin{example}
    There is a $1$-Lipschitz (but not convex), such that the above doesn't hold.
    On $\{0, 1\}^n$ take
    \[ f(x) = \sqrt{\max(-n/2 + \sum_i x_i, 0)} \]
    Alternatively,
    Let $A = \left\{(x_1, \dots, x_n) \in \{0, 1\}^n, \sum_i x_i \le \frac{n}{2}\right\}$
    and $f(x) = \inf_{y \in \{ (x_1, \dots, x_n) \in \{0, 1\}^n, \sum_i x_i \le \frac{n}{2} \}} \norm{x -y}_2$.
    Obviously this Lipschitz, as this is distance function. Furthermore,
    \begin{align*}
        \Pr[|f-\Exp f| > n^{1/4}] = \Pr[|\sqrt{\max\{ \frac{(2 x_i - 1) + \dots + (2 x_n - 1)}{n}, 0\}} - \Exp[f]| > \sqrt{2}]
    \end{align*}
    If the concentration held, the LHS would be $e^{-\sqrt{n}}$. But actually the thing inside the max looks like a Gaussian.
    This means we are after
    \[ \Pr[| \sqrt{\max\{\zeta, 0\}} - \Exp f| > \sqrt{2}] \]
    for some $\zeta \sim \N(0, I)$, but this is some constant.
\end{example}